\chapter{Tangential derivatives on overlapping pieces}\label{sec:tangential-differential-independence}
The tangential derivative of an ambient Lipschitz function is defined by restricting to smooth pieces. We prove that these definitions agree almost everywhere on overlaps. The source is \cite{Sim18}, Chapter 3, Section 2, printed p.~92, together with the chart differentiability discussion in Chapter 2, Section 4. We use the already proved tangent agreement and Rademacher theorem, and keep the area formula as a recorded input. Throughout, $U\subset\mathbb R^n$ is open and $1\leq m\leq n$.

\begin{lemma}[A density-one set approaches every direction]\label{lem:density-one-direction-approximation}
Let $A\subset\mathbb R^m$ be measurable and have Lebesgue density one at $a$. For every fixed $v\in\mathbb R^m$,
\[
 \operatorname{dist}(a+tv,A)=o(|t|)\quad\text{as }t\to0.
\]
Consequently, for any sequence $t_j\to0$ of nonzero numbers, there are $a_j\in A$ with $a_j=a+t_jv+o(|t_j|)$.
\end{lemma}
\begin{proof}
If the distance estimate failed, some $\varepsilon>0$ and sequence $t_j\to0$ would have $B_{\varepsilon|t_j|/2}(a+t_jv)$ disjoint from $A$. These balls lie in $B_{(|v|+\varepsilon)|t_j|}(a)$ and have volume $\omega_m(\varepsilon|t_j|/2)^m$. The fraction of missing volume in the latter balls would be bounded below by the positive number $(\varepsilon/(2(|v|+\varepsilon)))^m$, contradicting density one. This also covers $v=0$. To choose $a_j$, approximate the infimum defining the distance within $|t_j|/j$; density one ensures $A$ is nonempty near $a$. These points give the displayed expansion. The lemma provides points of the set approaching a direction, not a whole line segment contained in it.
\end{proof}

\begin{proposition}[Differentiable functions agreeing on a measurable set have the same derivative there]\label{prop:agreement-set-equal-derivatives}
Let $V\subset\mathbb R^m$ be open and $f,g:V\to\mathbb R^k$ be locally Lipschitz and let $A\subset V$ be measurable with $f=g$ on $A$. Then $Df=Dg$ at almost every point of $A$. More precisely this holds at every density-one point $a\in A$ where both functions are differentiable.
\end{proposition}
\begin{proof}
At such a point fix $v$ and use Lemma \ref{lem:density-one-direction-approximation} to choose $a_j\in A$ with $a_j=a+t_jv+o(|t_j|)$. Because $f(a)=g(a)$ and $f(a_j)=g(a_j)$, their differentiability expansions give
\[
 0=(Df(a)-Dg(a))(a_j-a)+o(|a_j-a|)
   =t_j(Df(a)-Dg(a))v+o(|t_j|).
\]
Divide by $t_j$ and pass to the limit. The difference vanishes on every $v$, so the derivatives agree. Rademacher and Lebesgue differentiation make the indicated differentiability and density conditions hold at almost every point of $A$. This argument applies componentwise to vector targets and does not require that $A$ have interior.
\end{proof}

\begin{lemma}[Restriction to a regular chart defines a tangential differential]\label{lem:chart-tangential-differential}
Let $F:V\subset\mathbb R^m\to N\subset U\subset\mathbb R^n$ be an injective regular $C^1$ chart of an embedded manifold, and let $f:U\to\mathbb R^k$ be locally Lipschitz. At almost every chart point $a$, the derivative $D(f\circ F)(a)$ exists and defines a linear map on $T_{F(a)}N$ by
\[
 L_{F(a)}(DF(a)v)=D(f\circ F)(a)v.
\]
On a neighborhood where $f$ has Lipschitz constant $L$, this map has operator norm at most $L$. Changing regular $C^1$ coordinates gives the same map at points where the derivatives exist.
\end{lemma}
\begin{proof}
On relatively compact base balls $F$ has bounded derivative, so $f\circ F$ is Lipschitz. Rademacher gives its derivative almost everywhere; a countable exhaustion covers the chart. Regularity makes $DF(a)$ an isomorphism onto its tangent image, giving a well-defined linear map. For each $v$, the Lipschitz bound for $f$ gives
\[
 |D(f\circ F)(a)v|
 =\lim_{t\to0}\frac{|f(F(a+tv))-f(F(a))|}{|t|}
 \leq L|DF(a)v|.
\]
The chain rule for composition with a $C^1$ coordinate diffeomorphism proves coordinate independence. Its derivative and inverse derivative are locally bounded, so null sets are preserved under such coordinate changes. The area formula, with the strictly positive chart Jacobian, identifies coordinate-null exceptional sets with $\mathcal H^m$-null exceptional subsets of $N$. An ambient derivative of $f$ need not exist at these surface points; only derivatives of its chart restriction are used.
\end{proof}

\begin{proposition}[Tangential differentials agree on two overlapping manifolds]\label{prop:overlap-tangential-differentials}
Let $N_1,N_2\subset U\subset\mathbb R^n$ be embedded $m$-dimensional $C^1$ manifolds and $f:U\to\mathbb R^k$ locally Lipschitz. At $\mathcal H^m$-almost every $x\in N_1\cap N_2$, their tangent planes coincide and the two tangential differentials of $f$ agree on that common plane. No open overlap of the manifolds is required.
\end{proposition}
\begin{proof}
Cover both manifolds by countably many finite-area graph patches and consider one pair of patches with charts $F_i$. The tangent agreement follows from Corollary \ref{cor:rectifiable-common-tangents}. The set $A=F_1^{-1}(N_1\cap N_2)$ is measurable in the first chart. Lebesgue differentiation makes almost every point of $A$ a density-one point. The positive Jacobian in the area formula transfers this conclusion to almost every point of the intersection. Lemma \ref{lem:chart-tangential-differential} removes further null sets so both chart restrictions are differentiable.

At a remaining point $x=F_1(a)$ and for $v\in T_xN_1=T_xN_2$, choose $w$ with $DF_1(a)w=v$. Lemma \ref{lem:density-one-direction-approximation} gives $a_j\in A$ with $a_j=a+t_jw+o(|t_j|)$. Put $x_j=F_1(a_j)$; then $x_j\in N_1\cap N_2$ and $x_j=x+t_jv+o(|t_j|)$. Use graph coordinates for the second patch. Its inverse is the restriction of a linear orthogonal projection $p_2$ onto its base plane, after translation. Thus its inverse coordinates satisfy $b_j=b+t_jp_2v+o(|t_j|)$, and $DF_2(b)p_2v=v$ because $v$ belongs to its tangent plane. Expand $f\circ F_i$ along these points. Both expansions evaluate the same quantity $f(x_j)-f(x)$, and they give respectively $t_jL_1(v)+o(|t_j|)$ and $t_jL_2(v)+o(|t_j|)$. Dividing and passing to the limit gives $L_1(v)=L_2(v)$. The finite-dimensional argument applies to each direction at the same good point. A countable union over the pairs of patches completes the proof.
\end{proof}

\begin{theorem}[An intrinsic tangential differential on a rectifiable set]\label{thm:rectifiable-tangential-differential}
Let $M\subset U$ be measurable, countably $m$-rectifiable and locally finite in $\mathcal H^m$ measure, and let $f:U\to\mathbb R^k$ be locally Lipschitz. Differentiation on any countable $C^1$ decomposition defines a measurable linear map
\[
 d^Mf_x:T_xM\to\mathbb R^k
 \quad\text{for }\mathcal H^m\text{-almost every }x\in M,
\]
independent of the decomposition up to a null set. Its norm is at most the local Lipschitz constant of $f$. Where the ambient derivative exists, $d^Mf_x=Df(x)|_{T_xM}$ at the almost-everywhere points under consideration. For scalar $f$, its Euclidean dual is the tangential gradient $\nabla^Mf$.
\end{theorem}
\begin{proof}
Use Theorem \ref{thm:rectifiable-c1-decomposition} and Lemma \ref{lem:finite-area-c1-patches} to partition $M$ into countably many measurable chart pieces modulo a null set. Lemma \ref{lem:chart-tangential-differential} defines the maps on each piece and bounds their norms. The derivative of a Lipschitz chart function is measurable as a limit of difference quotients. In a chart with $g=f\circ F$, the ambient representative is the matrix
\[
 Dg\bigl((DF)^*DF\bigr)^{-1}(DF)^*.
\]
It restricts to the required map on the tangent plane and is zero on its orthogonal complement. The inverse chart metric is continuous, so this matrix field is measurable. Thus the resulting field is measurable on the pieces and on their countable union.

For two decompositions, compare every pair of patches by Proposition \ref{prop:overlap-tangential-differentials}. The countable union of its exceptional sets is null. The maps therefore agree on one full-measure subset of $M$; tangent identification there is Theorem \ref{thm:rectifiable-forward-tangents}. If $f$ is ambient differentiable at such a point, its expansion along a $C^1$ chart immediately gives $D(f\circ F)=Df\circ DF$, proving the restriction formula. For a scalar target, the Euclidean inner product on the tangent plane identifies a linear functional with a unique vector. This proves the independence assertion used in Definition \ref{def:tangential-coarea-jacobian}; it does not prove the coarea formula or its almost-everywhere fiber geometry.
\end{proof}
